%=====================================================================
\chapter{Propositional Logic and Inference}\label{ch:propositional}
%=====================================================================
\locator{4}{Part IV --- Knowledge, Logic and Reasoning}

\begin{outcomes}
By the end of this chapter you will be able to:
\begin{tight}
  \item \textbf{Write} statements about a domain in propositional logic and
        \textbf{evaluate} them in a model.
  \item \textbf{Distinguish} entailment from derivation, and \textbf{explain} what
        soundness and completeness each guarantee.
  \item \textbf{Decide} entailment by model checking, by DPLL and by resolution
        refutation, and \textbf{compare} their cost.
  \item \textbf{Convert} a sentence to conjunctive normal form.
  \item \textbf{Produce} a countermodel when a conclusion does \emph{not} follow,
        and \textbf{explain} why that is the more useful answer in practice.
\end{tight}
\end{outcomes}

\begin{prereq}
\chref{kr} (what a representation is for, and the price defaults charge) and
\chref{python} (recursion and sets). No previous logic is assumed. \chref{csp} is
useful: deciding satisfiability is a constraint problem, and DPLL is backtracking
with propagation under another name.
\end{prereq}

\begin{keyterms}
proposition $\bullet$ connective $\bullet$ model $\bullet$ truth table $\bullet$
satisfiable $\bullet$ valid $\bullet$ entailment ($\models$) $\bullet$ derivation
($\vdash$) $\bullet$ soundness $\bullet$ completeness $\bullet$ literal $\bullet$
clause $\bullet$ conjunctive normal form $\bullet$ model checking $\bullet$ DPLL
$\bullet$ unit propagation $\bullet$ pure literal $\bullet$ resolution $\bullet$
refutation $\bullet$ countermodel
\end{keyterms}

\section{The Release Gate, Written Down}

The delivery team has a rule about when a release may go out. Everyone can state it
and nobody has written it down precisely, so every release involves an argument.
Here it is in logic.

\begin{examplebox}[The release gate]
Six propositions --- each either true or false --- and seven sentences.

\smallskip
\noindent\begin{tabular}{@{}L{0.9cm}L{6.0cm}L{7.0cm}@{}}
\toprule
 & \textbf{Sentence} & \textbf{In English} \\
\midrule
R1 & $\mathit{Reviewed} \wedge \mathit{Tested} \Rightarrow \mathit{Staged}$ &
Code that is reviewed and tested reaches staging \\
R2 & $\mathit{Staged} \wedge \mathit{SignedOff} \Rightarrow \mathit{Deploy}$ &
Staged and signed off means it may deploy \\
R3 & $\mathit{Blocker} \Rightarrow \neg \mathit{Deploy}$ &
An open blocker forbids deployment \\
F1 & $\mathit{Reviewed}$ & The code was reviewed \\
F2 & $\mathit{Tested}$ & The tests passed \\
F3 & $\mathit{SignedOff}$ & The product owner signed off \\
F4 & $\neg \mathit{Blocker}$ & There is no open blocker \\
\bottomrule
\end{tabular}

\smallskip
The question the team actually argues about: \emph{may we deploy?} In the notation
of this chapter, is $\mathit{KB} \models \mathit{Deploy}$?
\end{examplebox}

\section{Syntax, Semantics and Models}

\begin{definitionbox}[Model, satisfiability, validity]
A \term{model} is an assignment of true or false to every proposition symbol. With
$n$ symbols there are $2^{n}$ models; the gate above has six symbols and therefore
$64$.

A sentence is \term{satisfiable} if some model makes it true, \term{valid} if every
model does, and \term{unsatisfiable} if none does. The connectives have their usual
truth tables; the only one that regularly surprises people is implication:
\[
  \begin{array}{cc|c}
    a & b & a \Rightarrow b \\ \hline
    \text{F} & \text{F} & \text{T} \\
    \text{F} & \text{T} & \text{T} \\
    \text{T} & \text{F} & \text{F} \\
    \text{T} & \text{T} & \text{T}
  \end{array}
\]
$a \Rightarrow b$ is false in exactly one case: $a$ true and $b$ false. So
``if there is a blocker then we do not deploy'' says nothing whatever about the case
where there is no blocker --- which is correct, and is why F4 has to be stated
separately.
\end{definitionbox}

\begin{alertbox}[Implication is not causation, and not relevance]
$\mathit{Blocker} \Rightarrow \neg \mathit{Deploy}$ does not say the blocker
\emph{causes} the decision, and $a \Rightarrow b$ is true whenever $a$ is false
regardless of what $b$ is about. ``If the build number is prime then Ayesha is on
call'' is a true sentence on any day the build number is composite.

This bothers everyone at first and it is not a defect: it is the price of a
connective whose truth depends only on the truth of its parts, which is exactly the
property that makes mechanical inference possible.
\end{alertbox}

\section{Entailment and Derivation}

Two ideas that look alike, are constantly confused, and are the foundation of the
whole part.

\begin{definitionbox}[Entailment and derivation]
$\mathit{KB} \models \alpha$ --- \term{entailment} --- means $\alpha$ is true in
\emph{every} model in which $\mathit{KB}$ is true. It is a statement about
\emph{meaning}, and it makes no reference to any algorithm.

$\mathit{KB} \vdash_{i} \alpha$ --- \term{derivation} --- means the inference
procedure $i$ can produce $\alpha$ from $\mathit{KB}$. It is a statement about
\emph{that procedure}.

A procedure is \term{sound} if everything it derives is entailed
($\vdash$ implies $\models$): it never lies. It is \term{complete} if everything
entailed can be derived ($\models$ implies $\vdash$): it never misses anything.
\end{definitionbox}

An inference procedure is worth having exactly when it is both. Soundness without
completeness gives a system that is right when it speaks and sometimes silent;
completeness without soundness is worthless, because a procedure that derives
everything is complete.

\dsfig{%
\begin{tikzpicture}[font=\scriptsize]
\begin{scope}
  \fill[primary!6, rounded corners=4pt] (-4.6,-2.0) rectangle (4.6,2.0);
  \node[font=\scriptsize\bfseries, text=primary!70!black, anchor=north west]
    at (-4.5,1.92) {all 64 models};

  \draw[greenhl, line width=1.2pt, fill=greenhl!10] (1.2,0) ellipse (3.0 and 1.5);
  \node[font=\scriptsize\bfseries, text=greenhl!50!black] at (2.9,1.05)
    {models where \emph{Deploy} is true};

  \draw[primary, line width=1.4pt, fill=primary!22] (0.1,0) ellipse (1.25 and 0.78);
  \node[font=\scriptsize\bfseries, text=primary!70!black] at (0.1,0.28)
    {models of KB};
  \node[font=\tiny, text=primary!70!black] at (0.1,-0.12) {here: exactly 1};
\end{scope}

\node[font=\scriptsize, text=neutral, text width=4.2cm, align=flush left,
      anchor=north west] at (-4.5,-0.35)
  {\textbf{$\mathit{KB} \models \mathit{Deploy}$} means the inner region lies
   \emph{entirely} inside the outer one: there is no model in which the knowledge
   base holds and the conclusion fails.};

\node[rounded corners=3pt, fill=accent!7, draw=accent!55, line width=0.9pt,
      text=accent!55!black, font=\scriptsize, text width=13.0cm,
      align=flush left, inner sep=6pt] at (0,-2.85)
  {\textbf{More knowledge shrinks the inner region.} Every sentence added to the
   knowledge base rules models out, so the set of entailed conclusions can only
   grow. That is \term{monotonicity}, and it is exactly what \chref{kr}'s defaults
   gave up: there, adding ``Release 4.7 is a regulated change'' \emph{removed} a
   conclusion. Logic cannot do that, and in exchange it never has to guess.};
\end{tikzpicture}}%
{Entailment as containment between sets of models. The picture is the definition;
the algorithms in the rest of the chapter are three ways of checking
it.}{entailment}

\section{Deciding It: Model Checking}

The definition is directly executable: enumerate all $2^{n}$ models, keep those
satisfying every sentence of the knowledge base, and check whether the query is true
in all of them.

On the release gate this is decisive. Of $64$ models, exactly \textbf{one} satisfies
the knowledge base:
\[
  \mathit{Reviewed} = \mathit{Tested} = \mathit{SignedOff} = \mathit{Staged}
  = \mathit{Deploy} = \text{T}, \qquad \mathit{Blocker} = \text{F}
\]
and $\mathit{Deploy}$ is true in it. So $\mathit{KB} \models \mathit{Deploy}$: the
release may go out, and the argument is over.

Model checking is sound, complete, and trivially correct --- and it is $O(2^{n})$.
Six symbols is $64$; a realistic gate with forty conditions is $10^{12}$. It is the
definition made executable, not a usable algorithm.

\section{Deciding It Better: DPLL}

\term{DPLL} is a backtracking search over assignments --- exactly \chref{csp}'s
backtracking, with propagation --- plus two rules that do most of the work.

\begin{definitionbox}[DPLL]
Work on the knowledge base in \term{conjunctive normal form}: a conjunction of
\term{clauses}, each a disjunction of \term{literals}. To test satisfiability:
\begin{tight}
  \item If there are no clauses left, the assignment satisfies everything: report
        \emph{satisfiable}.
  \item If any clause is empty, it cannot be satisfied: backtrack.
  \item \term{Unit propagation}: if a clause has one literal left, that literal
        \emph{must} be true. Assign it and simplify.
  \item \term{Pure literal}: if a symbol appears with only one sign anywhere,
        assign it that way; it can never help to do otherwise.
  \item Otherwise pick a symbol and try both values.
\end{tight}
\end{definitionbox}

To decide entailment with a satisfiability test, use the equivalence that underlies
every refutation proof in this part:
\[
  \mathit{KB} \models \alpha
  \quad\text{if and only if}\quad
  (\mathit{KB} \wedge \neg\alpha) \text{ is unsatisfiable}
\]
If assuming the opposite of the conclusion leads to a contradiction, the conclusion
follows.

On the release gate, DPLL settles it in \textbf{6 recursive calls with 5 unit
propagations}, against model checking's $64$ enumerations --- and it never branches
once, because the facts F1--F4 are unit clauses that cascade.

\dsfig{%
\begin{tikzpicture}[font=\scriptsize,
  step/.style={rounded corners=3pt, draw=#1, line width=1pt, fill=#1!8,
               text=#1!50!black, text width=5.6cm, align=flush left,
               inner sep=5pt, font=\scriptsize},
  a/.style={-{Stealth[length=2mm]}, neutral!70, line width=0.9pt},
  tag/.style={rounded corners=2pt, fill=#1, text=white,
              font=\tiny\bfseries, inner sep=2.5pt}]

\node[step=primary] (s0) at (0,4.3)
  {\texttt{Reviewed} is a unit clause (F1). It must be true. Every clause
   containing it is satisfied and removed; \texttt{-Reviewed} is struck out of
   R1.};
\node[step=primary] (s1) at (0,2.9)
  {\texttt{Tested} (F2) is now a unit clause. R1 shortens to the single literal
   \texttt{Staged}.};
\node[step=greenhl] (s2) at (0,1.5)
  {\texttt{Staged} is therefore forced --- and it was never given as a fact. This
   is the first genuinely \emph{new} conclusion.};
\node[step=greenhl] (s3) at (0,0.1)
  {\texttt{SignedOff} (F3) plus \texttt{Staged} shortens R2 to \texttt{Deploy}.
   \texttt{Deploy} is forced.};
\node[step=accent] (s4) at (0,-1.3)
  {But we assumed \texttt{-Deploy} to test entailment. That clause is now
   \textbf{empty}: a contradiction.};

\foreach \a/\b in {s0/s1,s1/s2,s2/s3,s3/s4}{\draw[a] (\a) -- (\b);}

\node[tag=primary, anchor=west] at (3.1,4.3)  {unit};
\node[tag=primary, anchor=west] at (3.1,2.9)  {unit};
\node[tag=greenhl, anchor=west] at (3.1,1.5)  {derived};
\node[tag=greenhl, anchor=west] at (3.1,0.1)  {derived};
\node[tag=accent,  anchor=west] at (3.1,-1.3) {refutation};

\node[font=\scriptsize, text=neutral, text width=4.6cm, align=flush left,
      anchor=north west] at (3.9,3.9)
  {Five unit propagations, six recursive calls, and \emph{no branching at all}.

   \smallskip
   Model checking would have enumerated all 64 assignments to reach the same
   conclusion.

   \smallskip
   The contradiction means $\mathit{KB} \wedge \neg\mathit{Deploy}$ is
   unsatisfiable, and therefore
   $\mathit{KB} \models \mathit{Deploy}$.};
\end{tikzpicture}}%
{DPLL deciding the release gate. Unit propagation cascades from the four given
facts and closes the question without a single branch.}{dpll}

\section{Conjunctive Normal Form}

DPLL and resolution both need CNF. Every propositional sentence has an equivalent
one, obtained mechanically:

\begin{enumerate}[leftmargin=2.2em, itemsep=2pt, topsep=4pt]
  \item Eliminate $\Leftrightarrow$: replace $a \Leftrightarrow b$ with
        $(a \Rightarrow b) \wedge (b \Rightarrow a)$.
  \item Eliminate $\Rightarrow$: replace $a \Rightarrow b$ with
        $\neg a \vee b$.
  \item Drive $\neg$ inwards by De~Morgan:
        $\neg(a \wedge b) \equiv \neg a \vee \neg b$, and
        $\neg(a \vee b) \equiv \neg a \wedge \neg b$; remove double negations.
  \item Distribute $\vee$ over $\wedge$.
\end{enumerate}

Applying step~2 to the gate's rules gives the clauses the lab uses:

\smallskip
\noindent\begin{tabular}{@{}L{1.0cm}L{6.0cm}L{7.6cm}@{}}
\toprule
 & \textbf{Sentence} & \textbf{Clause} \\
\midrule
R1 & $\mathit{Reviewed} \wedge \mathit{Tested} \Rightarrow \mathit{Staged}$ &
$\neg \mathit{Reviewed} \vee \neg \mathit{Tested} \vee \mathit{Staged}$ \\
R2 & $\mathit{Staged} \wedge \mathit{SignedOff} \Rightarrow \mathit{Deploy}$ &
$\neg \mathit{Staged} \vee \neg \mathit{SignedOff} \vee \mathit{Deploy}$ \\
R3 & $\mathit{Blocker} \Rightarrow \neg \mathit{Deploy}$ &
$\neg \mathit{Blocker} \vee \neg \mathit{Deploy}$ \\
\bottomrule
\end{tabular}

\section{Resolution}

\begin{definitionbox}[The resolution rule]
From two clauses containing complementary literals, infer their combination with
that pair removed:
\[
  \frac{\alpha \vee \ell \qquad \beta \vee \neg\ell}
       {\alpha \vee \beta}
\]
Resolution is sound, and --- applied to CNF with the negated query added --- it is
\term{refutation complete}: if the set is unsatisfiable, repeated resolution will
derive the \term{empty clause}, which is a disjunction of nothing and therefore
false.
\end{definitionbox}

The proof for the release gate is five steps. Set as a table, because a proof is a
numbered list of clauses with a provenance column:

\smallskip
\rowcolors{2}{white}{lightbg}
\noindent\begin{longtable}{@{}L{1.0cm}L{6.6cm}L{7.0cm}@{}}
\toprule
\rowcolor{primary!15}
\textbf{\#} & \textbf{Clause} & \textbf{Derived from} \\
\midrule
\endhead
1 & $\mathit{Reviewed}$ & given (F1) \\
2 & $\mathit{Tested}$ & given (F2) \\
3 & $\mathit{SignedOff}$ & given (F3) \\
4 & $\neg \mathit{Reviewed} \vee \neg \mathit{Tested} \vee \mathit{Staged}$ &
given (R1) \\
5 & $\neg \mathit{Staged} \vee \neg \mathit{SignedOff} \vee \mathit{Deploy}$ &
given (R2) \\
6 & $\neg \mathit{Deploy}$ & \textbf{assumed}: the negation of the query \\
7 & $\neg \mathit{Tested} \vee \mathit{Staged}$ & 1, 4 \\
8 & $\mathit{Staged}$ & 2, 7 \\
9 & $\neg \mathit{Staged} \vee \mathit{Deploy}$ & 3, 5 \\
10 & $\mathit{Deploy}$ & 8, 9 \\
11 & $\square$ \quad (the empty clause) & 6, 10 \\
\bottomrule
\caption{A resolution refutation of $\neg \mathit{Deploy}$. Deriving the empty
clause shows the assumption is inconsistent with the knowledge base, so
$\mathit{KB} \models \mathit{Deploy}$.}\label{tab:resolution}
\end{longtable}

\begin{alertbox}[Resolution is not the fast one]
That proof is five steps, and a naive implementation does not find it in five. The
lab's breadth-first version --- resolve every pair, add everything new, repeat ---
takes \textbf{105 resolution steps} to reach the empty clause, which is more work
than enumerating all $64$ models.

So resolution is not here because it is efficient on propositional problems; DPLL
is, and modern SAT solvers are DPLL with better bookkeeping. Resolution is here for
two reasons. It produces an \emph{auditable proof} --- Table~\ref{tab:resolution}
can be read by a person and checked line by line, which no model count can. And it
\emph{lifts to first-order logic} (\chref{fol}), where there are infinitely many
models and enumerating them is not merely slow but impossible.
\end{alertbox}

\section{When the Answer Is No}

Drop F3 --- suppose the product owner has not signed off. Now three models satisfy
the knowledge base, and $\mathit{Deploy}$ is false in one of them:
\[
  \mathit{Reviewed} = \mathit{Tested} = \mathit{Staged} = \text{T},\quad
  \mathit{SignedOff} = \mathit{Blocker} = \mathit{Deploy} = \text{F}
\]
That is a \term{countermodel}: a concrete situation consistent with everything known
in which the conclusion fails. Resolution correctly fails to derive the empty
clause, and model checking hands you the countermodel directly.

In practice the countermodel is the more valuable output. ``It does not follow'' is
an assertion a team can argue with; ``here is a situation where everything you told
me holds and the release still must not go out --- nobody signed off'' ends the
argument. When you build one of these systems, return the countermodel.

\begin{worked}[A conclusion the knowledge base implies without being told]
Drop F4 instead --- delete $\neg \mathit{Blocker}$, so nobody has checked the
blocker list. Does $\mathit{Deploy}$ still follow?

\smallskip
\textbf{The intuition says no.} Nothing now rules out a blocker, and R3 says a
blocker forbids deployment.

\smallskip
\textbf{The models say yes.} Work it through:
\begin{tight}
  \item F1, F2, F3 force $\mathit{Reviewed}$, $\mathit{Tested}$,
        $\mathit{SignedOff}$ all true.
  \item R1 then forces $\mathit{Staged}$ true.
  \item R2 then forces $\mathit{Deploy}$ true.
  \item R3 says $\mathit{Blocker} \Rightarrow \neg\mathit{Deploy}$. Since
        $\mathit{Deploy}$ is true, $\neg\mathit{Deploy}$ is false, so
        $\mathit{Blocker}$ \emph{must be false} --- otherwise R3 is violated.
\end{tight}
Exactly one model survives, and it is the same one as before. So
$\mathit{KB} \models \mathit{Deploy}$, \emph{and} $\mathit{KB} \models \neg
\mathit{Blocker}$: the knowledge base entails that there is no blocker, without
anyone having said so.

\smallskip
\textbf{Is that right?} Logically, yes, and it is worth sitting with the
discomfort. The knowledge base asserts as facts that the code is reviewed, tested
and signed off, and asserts as a rule that these permit deployment. Taken together
those \emph{are inconsistent with} an open blocker. The system has not been
careless; it has drawn a consequence of what it was told.

\smallskip
\textbf{And it is a modelling error.} The team's intent was surely ``deployment is
permitted \emph{provided} no blocker'', which is
$\mathit{Staged} \wedge \mathit{SignedOff} \wedge \neg \mathit{Blocker}
\Rightarrow \mathit{Deploy}$. Write R2 that way and the entailment disappears:
without F4, $\mathit{Deploy}$ no longer follows, which is what everyone expected.

\smallskip
\textbf{The lesson.} Inference is only as good as the axioms, and the failure mode
is not a wrong answer --- it is a \emph{correct answer to a question you did not
mean to ask}. Contrast \chref{kr}: a default would have said ``normally no blocker''
and withdrawn it on evidence. Logic will not guess, so everything you mean must be
written down, including the things that seem too obvious to state.
\end{worked}

\begin{pitfall}
\begin{tight}
  \item \textbf{Confusing $\models$ with $\vdash$.} The first is about meaning, the
        second about a procedure. Soundness and completeness are precisely the two
        claims that connect them.
  \item \textbf{Expecting $a \Rightarrow b$ to mean $a$ causes $b$.} It is false in
        exactly one row of the truth table and is true whenever $a$ is false.
  \item \textbf{Omitting a condition from a rule because it is obvious.} The worked
        example above turns on exactly this, and the result is a system that
        concludes there is no blocker.
  \item \textbf{Using model checking beyond about twenty symbols.} $2^{n}$ is the
        definition, not an algorithm.
  \item \textbf{Assuming resolution is the efficient method.} It took $105$ steps
        where DPLL took $6$ calls. Its value is auditable proofs and the lift to
        first-order logic.
  \item \textbf{Reporting only ``does not follow''.} Return the countermodel; it is
        what turns a refusal into a diagnosis.
\end{tight}
\end{pitfall}

%---------------------------------------------------------------------
\begin{lab}[Three Ways to Decide a Release Gate]
Model checking, DPLL and resolution refutation on the same knowledge base, plus the
two variants of the worked example.
\begin{lstlisting}[language=Python]
"""Chapter 11 lab -- deciding a release gate three ways.

A clause is a frozenset of literals; a literal is "X" or "-X". The
three procedures agree on every question, and differ enormously in
what they cost and in what they hand back.
"""

import itertools

SYMBOLS = ["Reviewed", "Tested", "SignedOff", "Blocker", "Staged",
           "Deploy"]

# The knowledge base, already in conjunctive normal form.
KB = [
    frozenset({"-Reviewed", "-Tested", "Staged"}),      # R1
    frozenset({"-Staged", "-SignedOff", "Deploy"}),     # R2
    frozenset({"-Blocker", "-Deploy"}),                 # R3
    frozenset({"Reviewed"}),                            # F1
    frozenset({"Tested"}),                              # F2
    frozenset({"SignedOff"}),                           # F3
    frozenset({"-Blocker"}),                            # F4
]
LABEL = {KB[0]: "R1", KB[1]: "R2", KB[2]: "R3", KB[3]: "F1",
         KB[4]: "F2", KB[5]: "F3", KB[6]: "F4"}


def neg(lit):
    return lit[1:] if lit.startswith("-") else "-" + lit


def holds(lit, model):
    return (not model[lit[1:]]) if lit.startswith("-") else model[lit]


def clause_true(clause, model):
    return any(holds(l, model) for l in clause)


# --- 1. model checking: the definition, made executable ------------
def model_check(kb, query, symbols):
    """Returns (entailed, models of kb, models enumerated)."""
    models, seen = [], 0
    for combo in itertools.product([False, True], repeat=len(symbols)):
        m = dict(zip(symbols, combo))
        seen += 1
        if all(clause_true(c, m) for c in kb):
            models.append(m)
    entailed = bool(models) and all(m[query] for m in models)
    return entailed, models, seen


# --- 2. DPLL: backtracking with propagation ------------------------
class Work:
    def __init__(self):
        self.calls = 0
        self.units = 0
        self.branches = 0


def simplify(kb, lit):
    out = []
    for c in kb:
        if lit in c:
            continue                       # clause already satisfied
        if neg(lit) in c:
            out.append(c - {neg(lit)})     # clause shortens
        else:
            out.append(c)
    return out


def dpll(kb, assign, w):
    w.calls += 1
    if not kb:
        return True, assign
    if any(len(c) == 0 for c in kb):
        return False, None
    for c in kb:                           # unit propagation
        if len(c) == 1:
            lit = next(iter(c))
            w.units += 1
            a = dict(assign)
            a[lit.lstrip("-")] = not lit.startswith("-")
            return dpll(simplify(kb, lit), a, w)
    lits = {l for c in kb for l in c}
    for l in lits:                         # pure literal
        if neg(l) not in lits:
            a = dict(assign)
            a[l.lstrip("-")] = not l.startswith("-")
            return dpll(simplify(kb, l), a, w)
    l = next(iter(next(iter(kb))))         # branch
    w.branches += 1
    for choice in (l, neg(l)):
        a = dict(assign)
        a[choice.lstrip("-")] = not choice.startswith("-")
        ok, res = dpll(simplify(kb, choice), a, w)
        if ok:
            return True, res
    return False, None


def entails_dpll(kb, query):
    """KB |= q  iff  KB and not-q is unsatisfiable."""
    w = Work()
    sat, _ = dpll(list(kb) + [frozenset({neg(query)})], {}, w)
    return (not sat), w


# --- 3. resolution refutation --------------------------------------
def resolve(ci, cj):
    out = []
    for l in ci:
        if neg(l) in cj:
            out.append((ci - {l}) | (cj - {neg(l)}))
    return out


def resolution(kb, query):
    """Breadth-first saturation. Honest, complete, and slow."""
    clauses = list(kb) + [frozenset({neg(query)})]
    steps = 0
    while True:
        new = set()
        for i in range(len(clauses)):
            for j in range(i + 1, len(clauses)):
                for r in resolve(clauses[i], clauses[j]):
                    steps += 1
                    if not r:
                        return True, steps
                    if r not in clauses:
                        new.add(r)
        if not new - set(clauses):
            return False, steps
        clauses += sorted(new, key=lambda c: (len(c), sorted(c)))


def show(model):
    return "  ".join("%s=%d" % (s, model[s]) for s in SYMBOLS)


if __name__ == "__main__":
    print("symbols %d, so %d models to enumerate"
          % (len(SYMBOLS), 2 ** len(SYMBOLS)))
    print()

    ent, models, seen = model_check(KB, "Deploy", SYMBOLS)
    print("MODEL CHECKING")
    print("   %d of %d models satisfy the KB" % (len(models), seen))
    for m in models:
        print("      " + show(m))
    print("   KB entails Deploy:", ent)
    assert ent and len(models) == 1

    ent2, w = entails_dpll(KB, "Deploy")
    print()
    print("DPLL")
    print("   %d calls, %d unit propagations, %d branches"
          % (w.calls, w.units, w.branches))
    print("   KB entails Deploy:", ent2)
    assert ent2 and ent2 == ent
    assert w.branches == 0, "the facts cascade; no branching needed"

    ent3, steps = resolution(KB, "Deploy")
    print()
    print("RESOLUTION")
    print("   empty clause derived:", ent3, " in %d steps" % steps)
    assert ent3 == ent
    print("   ... which is more work than enumerating all %d models."
          % seen)
    assert steps > seen

    # --- drop the sign-off: the answer becomes no ------------------
    print()
    print("=== drop F3 (SignedOff) ===")
    KB2 = [c for c in KB if c != frozenset({"SignedOff"})]
    ent4, models4, _ = model_check(KB2, "Deploy", SYMBOLS)
    print("   %d models satisfy the KB; entails Deploy: %s"
          % (len(models4), ent4))
    counter = [m for m in models4 if not m["Deploy"]]
    print("   countermodel:")
    print("      " + show(counter[0]))
    assert not ent4 and counter
    ok4, _ = resolution(KB2, "Deploy")
    assert not ok4
    print("   resolution correctly fails to derive the empty clause.")
    print("   The countermodel is the useful output: everything you")
    print("   told me holds, and the release still must not go out.")

    # --- drop the blocker fact: the surprise ----------------------
    print()
    print("=== drop F4 (not Blocker) ===")
    KB3 = [c for c in KB if c != frozenset({"-Blocker"})]
    ent5, models5, _ = model_check(KB3, "Deploy", SYMBOLS)
    print("   %d models satisfy the KB; entails Deploy: %s"
          % (len(models5), ent5))
    assert ent5 and len(models5) == 1
    print("   Blocker is %d in the only model: the KB entails that"
          % models5[0]["Blocker"])
    print("   there is no blocker, without anyone saying so.")
    assert not models5[0]["Blocker"]

    # --- and the fix: put the condition in the rule ---------------
    print()
    print("=== the fix: R2 requires no blocker ===")
    KB4 = [c for c in KB3
           if c != frozenset({"-Staged", "-SignedOff", "Deploy"})]
    KB4.append(frozenset({"-Staged", "-SignedOff", "Blocker", "Deploy"}))
    ent6, models6, _ = model_check(KB4, "Deploy", SYMBOLS)
    print("   %d models satisfy the KB; entails Deploy: %s"
          % (len(models6), ent6))
    assert not ent6
    print("   Now Deploy does NOT follow without knowing about")
    print("   blockers -- which is what the team meant all along.")
\end{lstlisting}
Then add a fortieth condition to the gate and re-run. Model checking stops being
possible ($2^{40}$ models); DPLL does not notice. That is the moment the difference
between a definition and an algorithm becomes concrete.
\end{lab}

%---------------------------------------------------------------------
\begin{checkpoint}
\begin{tightnum}
  \item State the difference between $\mathit{KB} \models \alpha$ and
        $\mathit{KB} \vdash \alpha$, and say which of soundness and completeness
        would be violated by a procedure that derives every sentence.
  \item With F4 removed, the knowledge base still entails $\mathit{Deploy}$
        \emph{and} entails $\neg \mathit{Blocker}$. Explain how, and say whether
        this is a bug in the inference procedure.
  \item Resolution took $105$ steps where DPLL took $6$ calls on the same problem.
        Give the two reasons resolution is nevertheless worth teaching.
\end{tightnum}
\end{checkpoint}

%---------------------------------------------------------------------
\begin{chaptersummary}
\begin{tight}
  \item A \term{model} assigns true or false to every symbol; $n$ symbols give
        $2^{n}$ models. A sentence is \term{satisfiable} if some model makes it
        true and \term{valid} if all do.
  \item $\mathit{KB} \models \alpha$ means $\alpha$ holds in every model of
        $\mathit{KB}$ --- a statement about meaning. $\mathit{KB} \vdash \alpha$
        means a procedure derives it. \term{Soundness} and \term{completeness} are
        the two claims linking them.
  \item Logic is \term{monotonic}: adding knowledge only ever rules models out, so
        conclusions are never withdrawn. That is precisely what \chref{kr}'s
        defaults traded away.
  \item \term{Model checking} is the definition made executable and is $O(2^{n})$.
        \term{DPLL} is backtracking with \term{unit propagation} and the
        \term{pure literal} rule --- $6$ calls against $64$ enumerations here, with
        no branching at all.
  \item $\mathit{KB} \models \alpha$ if and only if $\mathit{KB} \wedge \neg\alpha$
        is unsatisfiable. Every refutation method rests on this.
  \item \term{Resolution} on CNF is refutation complete. It is slower than DPLL
        ($105$ steps against $6$ calls) and is worth having because it yields an
        auditable proof and because it lifts to first-order logic.
  \item When a conclusion does not follow, return the \term{countermodel}. And
        remember that inference is only as good as the axioms: an omitted condition
        produces a correct answer to a question you did not mean to ask.
\end{tight}
\end{chaptersummary}

\begin{reviewq}
  \item \textbf{(MCQ)} $\mathit{KB} \models \alpha$ means: (a)~a procedure derives
        $\alpha$ (b)~$\alpha$ is true in every model of $\mathit{KB}$
        (c)~$\alpha$ is valid (d)~$\mathit{KB}$ is satisfiable.
  \item \textbf{(MCQ)} A sound inference procedure: (a)~derives everything entailed
        (b)~derives only things that are entailed (c)~always terminates
        (d)~is complete.
  \item \textbf{(MCQ)} $\mathit{KB} \models \alpha$ if and only if:
        (a)~$\mathit{KB} \wedge \alpha$ is satisfiable
        (b)~$\mathit{KB} \wedge \neg\alpha$ is unsatisfiable
        (c)~$\mathit{KB} \vee \alpha$ is valid (d)~$\alpha$ is a unit clause.
  \item \textbf{(MCQ)} Unit propagation applies when a clause has:
        (a)~no literals (b)~exactly one literal left (c)~only positive literals
        (d)~duplicate literals.
  \item \textbf{(Short)} Convert
        $\mathit{Staged} \wedge \mathit{SignedOff} \wedge \neg\mathit{Blocker}
        \Rightarrow \mathit{Deploy}$ to a clause, showing each step.
  \item \textbf{(Short)} Why is a countermodel a more useful answer than ``it does
        not follow''? Give the countermodel from this chapter and the sentence you
        would say to the delivery team.
  \item \textbf{(Short)} Resolution is slower than DPLL on propositional problems.
        Give the two reasons it is still the method developed in \chref{fol}.
  \item \textbf{(Applied)} Add the rule ``a release with an open blocker may still
        deploy if the incident commander approves'' to the gate, introducing one
        new symbol. Write it as a sentence, convert it to a clause, and determine
        --- by hand, using unit propagation --- whether $\mathit{Deploy}$ follows
        from F1--F3 together with $\mathit{Blocker}$ and the commander's approval.
\end{reviewq}
